%% [AI DRAFT] Borrador generado para discusión (Modo C, etapa 3).
%% El autor debe releerlo y reexpresarlo en su propia voz antes de darlo por bueno.
%% Todas las citas fueron verificadas contra el PDF; no hay marcadores [VERIFY] pendientes.

\section{MÉTODO}

Esta sección describe la arquitectura propuesta. La unidad de análisis es una \emph{ficha de
producto} que el comprador tiene abierta en su navegador: el sistema no recorre catálogos ni
consulta interfaces privilegiadas de la plataforma, sino que opera sobre la misma información
que el comprador ya tiene delante. Esa restricción gobierna todas las decisiones de diseño
que siguen y separa esta propuesta de los sistemas de detección desplegados por el operador
del comercio \cite{ref28}.

La arquitectura se organiza en tres capas (cliente, motor de riesgo y persistencia) y se
detalla en la Fig.~\ref{fig:arquitectura}.

\begin{figure*}[t]
\centering
\resizebox{0.95\textwidth}{!}{% ===================================================================
%  Fig. 1 — Arquitectura logica propuesta
%  Entregable: short-paper-es
%  Figura vectorial en TikZ: se edita como texto y compila con el paper.
%  Preambulo requerido:
%    \usepackage{tikz}
%    \usetikzlibrary{arrows.meta,positioning,fit,backgrounds,shapes.geometric}
%  Va dentro de un entorno figure* (ancho de las dos columnas).
% ===================================================================
\begin{tikzpicture}[
  font=\footnotesize,
  % Estilos base
  banda/.style   = {draw=black!40, rounded corners=4pt, line width=0.8pt},
  caja/.style    = {draw=black!60, fill=white, rounded corners=3pt,
                    align=center, minimum height=10mm, inner sep=4pt,
                    line width=0.6pt},
  cliente/.style = {caja, fill=orange!15, text width=25mm},
  servicio/.style= {caja, fill=blue!12, text width=28mm},
  persist/.style = {caja, fill=cyan!10, text width=30mm},
  api/.style     = {caja, fill=blue!20, text width=35mm, minimum height=14mm},
  etiq/.style    = {font=\scriptsize\bfseries, black!70},
  fl/.style      = {-{Latex[length=2.5mm]}, line width=0.7pt, black!70},
  flInv/.style   = {Latex[length=2.5mm]-, line width=0.7pt, black!70},
  flBid/.style   = {Latex[length=2.5mm]-Latex[length=2.5mm], line width=0.7pt, black!70},
  nota/.style    = {font=\tiny\itshape, black!50, align=center, fill=white,
                    inner sep=2pt, rounded corners=2pt},
]

% =============== CAPA 1: CLIENTE (arriba) ===============
\begin{scope}[on background layer]
  \node[banda, fill=gray!5, minimum width=160mm, minimum height=40mm,
        rounded corners=5pt] at (8,8) {};
\end{scope}
\node[etiq] at (8,9.8) {\textbf{Cliente}};

% Nodos del cliente
\node[cliente, fill=orange!25, text width=28mm, minimum height=15mm] (ficha) at (2.5,8)
  {\textbf{Ficha de}\\[-1pt]\textbf{Producto}\\[-3pt]
   \tiny\textbf{DOM}};

\node[cliente, fill=green!15, text width=30mm, minimum height=15mm] (extn) at (8,8)
  {\textbf{Extensión de}\\[-1pt]\textbf{navegador}};

\node[cliente, fill=cyan!15, text width=28mm, minimum height=15mm] (app) at (13.5,8)
  {\textbf{Aplicación}\\[-1pt]\textbf{web}};

% Flechas cliente
\draw[fl] (ficha) -- node[nota, above, yshift=1mm] {Envía información} (extn);
\draw[flBid] (extn) -- node[nota, above, yshift=1mm] {Solicita autenticación} (app);

% =============== BANDA DE CONEXIÓN ===============
\draw[fl] (extn.south) -- ++(0,-1.2) -| node[nota, above, pos=0.2] {Envía información extraída de la ficha de producto} (8,5.8);

% =============== CAPA 2: MOTOR DE RIESGO (medio) ===============
\begin{scope}[on background layer]
  \node[banda, fill=gray!5, minimum width=160mm, minimum height=60mm,
        rounded corners=5pt] at (8,3.5) {};
\end{scope}
\node[etiq] at (8,6.2) {\textbf{Motor de Riesgo}};

% API centralizadora
\node[api, fill=blue!25, minimum height=16mm] (api) at (8,4.5)
  {\textbf{API centralizadora de estimación}\\[-1pt]
   \textbf{del motor de riesgo}};

% Servicios del motor de riesgo
\node[servicio, fill=blue!15, text width=25mm, minimum height=12mm] (dom) at (1.5,2)
  {\textbf{Analizador de veracidad}\\[-1pt]\textbf{de dominio}};

\node[servicio, fill=blue!15, text width=25mm, minimum height=12mm] (rep) at (5.5,2)
  {\textbf{Analizador de}\\[-1pt]\textbf{reputación}};

\node[servicio, fill=blue!15, text width=28mm, minimum height=12mm] (rf) at (10.5,2)
  {\textbf{API de estimación del nivel}\\[-1pt]\textbf{de probabilidad de riesgo}};

\node[servicio, fill=blue!15, text width=28mm, minimum height=12mm] (bert) at (14.5,2)
  {\textbf{API de Detección de}\\[-1pt]\textbf{reseñas falsas}};

% Flechas de API a servicios
\draw[fl] (api.south) -- ++(0,-0.5) -| node[nota, above, pos=0.3] {Solicita análisis de veracidad de dominio} (dom.north);
\draw[fl] (api.south) -- ++(0,-0.5) -| node[nota, above, pos=0.3] {Solicita análisis de reputación del vendedor} (rep.north);
\draw[fl] (api.south) -- ++(0,-0.5) -| node[nota, above, pos=0.3] {Solicita análisis del nivel de probabilidad de riesgo} (rf.north);
\draw[fl] (api.south) -- ++(0,-0.5) -| node[nota, above, pos=0.3] {Solicita porcentaje de reseñas falsas} (bert.north);

% =============== BANDA DE CONEXIÓN A PERSISTENCIA ===============
\draw[fl] (api.south) -- ++(0,-0.8) -| node[nota, above, pos=0.2] {Almacena información de consultas realizadas} (8,-0.8);

% =============== CAPA 3: PERSISTENCIA (abajo) ===============
\begin{scope}[on background layer]
  \node[banda, fill=gray!5, minimum width=160mm, minimum height=35mm,
        rounded corners=5pt] at (8,-2.5) {};
\end{scope}
\node[etiq] at (8,-0.8) {\textbf{Persistencia}};

% Base de datos
\node[persist, fill=cyan!15, text width=35mm, minimum height=18mm] (bd) at (8,-3)
  {\textbf{Base de Datos}\\[-3pt]
   \tiny Consultas · características · veredictos\\[-1pt]
   \tiny (trazabilidad y reentrenamiento)};

% Flecha a base de datos
\draw[fl] (8,-1.5) -- (bd.north);

\end{tikzpicture}
}
\caption{Arquitectura lógica propuesta. La extensión lee el DOM ya renderizado de una página
abierta por el propio usuario; el motor de riesgo enriquece esas observaciones con análisis de
dominio e inferencia sobre vendedor y reseñas, y el agregador emite el veredicto con su nivel
de confianza.}
\label{fig:arquitectura}
\end{figure*}

\subsection{Capa cliente: observación sin recolección automatizada}

La capa cliente se implementa como una \emph{aplicación nativa en el navegador} (una
extensión) que se activa cuando el usuario abre una ficha de producto en una plataforma
soportada. La extensión recorre el árbol DOM ya
renderizado, localiza mediante selectores los elementos que exponen los atributos de interés
(precio, puntuación, número de reseñas, identidad del vendedor, medios de pago anunciados,
políticas de garantía) y construye un vector de observaciones que envía al motor de riesgo.
La extensión constituye el artefacto principal de esta propuesta, y no por conveniencia de
implementación. Es el único punto del sistema que coincide, en espacio y en tiempo, con el
instante en que el comprador decide: se ejecuta sobre la ficha que él tiene abierta, sin
mediación de la plataforma y antes de que el pago ocurra. Ningún otro componente de esta
arquitectura, ni las aproximaciones revisadas en la Sección~II, coincide de ese modo con el
instante de la decisión.

Un cliente complementario consume el mismo servicio y cubre la compra que ocurre fuera del
navegador de escritorio. Su papel es de cobertura, no de sustitución: opera sobre lo que el
usuario le proporciona o la plataforma expone públicamente, sin acceso al documento
renderizado, de modo que resolverá previsiblemente menos características. Su forma concreta,
web o móvil, no la determina la arquitectura, que solo exige consumir la misma interfaz y
respetar la misma regla de abstención.

Esta capa existe porque, según la evidencia recogida en la Sección~I, la advertencia que ha
de llegar antes del pago debe producirla el sistema y no el juicio del comprador
\cite{ref27, ref3}. La ventaja técnica de resolverlo en el cliente es que la información
necesaria ya fue entregada al navegador: no hay navegación automatizada ni recorrido de
múltiples páginas, que es lo que define al \textit{web scraping} \cite{ref24}, y el sistema no
emite hacia la plataforma del vendedor ninguna petición que el comprador no hubiera cursado en
cualquier caso.

«Del lado del cliente» no significa, sin embargo, «sin servidor»: el vector observado viaja al
motor de riesgo, el analizador de dominio consulta servicios de registro y de resolución de
nombres, y cada consulta queda registrada. La propuesta elimina la dependencia respecto de la
plataforma del vendedor a costa de introducir un nuevo tercero de confianza, con la política
de retención y anonimización que ello exige. Su otra contrapartida, declarada en la
Sección~V-C, es la fragilidad ante cambios de maquetación, pues los selectores predefinidos
son sensibles a modificaciones estructurales del sitio \cite{ref24}.

\begin{figure*}[t]
\centering
\resizebox{0.95\textwidth}{!}{% ===================================================================
%  Fig. 2 — Diagrama de secuencia: de la apertura de la ficha al veredicto
%  Entregable: short-paper-es
%  Preambulo requerido: tikz + arrows.meta, positioning, fit, backgrounds
%  Va dentro de un entorno figure* (ancho de las dos columnas).
% ===================================================================
\begin{tikzpicture}[
  font=\footnotesize, x=1cm, y=1cm,
  actor/.style   = {draw=black!75, fill=black!5, rounded corners=1.5pt,
                    align=center, minimum height=7.5mm, inner xsep=4pt,
                    text width=22mm, line width=0.4pt},
  vida/.style    = {densely dashed, black!45, line width=0.35pt},
  msg/.style     = {-{Latex[length=1.7mm]}, line width=0.45pt, black!80},
  ret/.style     = {-{Latex[length=1.7mm]}, line width=0.45pt, black!55,
                    densely dashed},
  et/.style      = {font=\scriptsize, inner sep=1.6pt, fill=white, align=center},
  nota/.style    = {font=\scriptsize\itshape, align=left, black!65,
                    draw=black!35, fill=black!3, rounded corners=1.5pt,
                    inner sep=4pt, text width=30mm},
]
\def\W{0.15}                      % media anchura de la barra de activación
\newcommand{\activa}[3]{%         % #1 = x, #2 = y inicio, #3 = y fin
  \fill[black!10] (#1-\W,#2) rectangle (#1+\W,#3);
  \draw[black!70, line width=0.35pt] (#1-\W,#2) rectangle (#1+\W,#3);}

% ---------- actores y líneas de vida ----------
\foreach \x/\n/\t in {
  0/usr/{Comprador},
  3.6/ext/{\textbf{Extensión}\\\scriptsize artefacto principal},
  7.6/dom/{Analizador\\de dominio},
  11.0/mod/{Modelos de\\inferencia},
  14.3/agg/{Agregador},
  17.3/bd/{BD de\\registros}} {
  \node[actor] (\n) at (\x,0) {\t};
  \draw[vida] (\x,-0.46) -- (\x,-10.5);
}

% ---------- activaciones ----------
\activa{3.6}{-1.2}{-10.0}
\activa{7.6}{-3.4}{-4.3}
\activa{11.0}{-5.1}{-6.0}
\activa{14.3}{-6.8}{-8.8}
\activa{17.3}{-8.4}{-8.8}

% ---------- mensajes ----------
\draw[msg] (0,-1.2) -- node[et,above]{abre la ficha de producto} (3.45,-1.2);

\draw[msg] (3.75,-2.0) -- ++(1.3,0) -- ++(0,-0.55)
   node[et,right,xshift=1.2mm,yshift=2mm]{lee el DOM ya renderizado} -- (3.75,-2.55);

\draw[msg] (3.75,-3.4) -- node[et,above]{nombre de dominio} (7.45,-3.4);
\draw[ret] (7.45,-4.3) -- node[et,above]{veracidad estimada} (3.75,-4.3);

\draw[msg] (3.75,-5.1) -- node[et,above]{características extraídas} (10.85,-5.1);
\draw[ret] (10.85,-6.0) -- node[et,above]{riesgo del vendedor · reseñas falsas} (3.75,-6.0);

\draw[msg] (3.75,-6.8) -- node[et,above]{resultados parciales} (14.15,-6.8);

\draw[msg] (14.45,-7.5) -- ++(1.15,0) -- ++(0,-0.55)
   node[et,right,xshift=1.2mm,yshift=2mm]{aplica la abstención} -- (14.45,-8.05);

\draw[msg] (14.45,-8.4) -- node[et,above]{registro de la consulta} (17.15,-8.4);

\draw[ret] (14.15,-9.3) -- node[et,above]{veredicto + nivel de confianza} (3.75,-9.3);
\draw[ret] (3.45,-10.0) -- node[et,above]{advertencia explicada, antes del pago} (0,-10.0);

% ---------- nota ----------
\node[nota, anchor=north west] at (18.2,-1.0)
  {La extensión no consulta a la plataforma del vendedor: opera sobre el documento que el navegador ya recibió.};

\end{tikzpicture}
}
\caption{Secuencia de una consulta, desde la apertura de la ficha hasta la advertencia. Las
barras verticales indican los intervalos de actividad de cada componente; las flechas
discontinuas, los retornos. Ninguna flecha sale hacia la plataforma del vendedor: las únicas
comunicaciones salientes van de la extensión a los componentes del motor de riesgo, y las
consultas del analizador a los servicios de registro y de resolución de nombres se resumen en
su intervalo de actividad. El agregador aplica la regla de abstención antes de emitir el
veredicto.}
\label{fig:secuencia}
\end{figure*}

La Fig.~\ref{fig:secuencia} traza el orden temporal de una consulta y precisa lo que la vista
estática no expresa: entre la apertura de la ficha y la advertencia no se emite ninguna
petición hacia la plataforma del vendedor.

\subsection{Conjunto de características y regla de abstención ante evidencia ausente}

El motor de riesgo trabaja sobre ocho características, resumidas en la Tabla~\ref{tab:senales}. Cada
característica se define junto con la capa que la computa, su dominio de valores y (de manera
central para esta propuesta) el valor que toma cuando no puede determinarse.

\begin{table}[t]
\caption{Conjunto de características, origen y comportamiento ante evidencia ausente}
\label{tab:senales}
\centering
\footnotesize
\setlength{\tabcolsep}{3pt}
\renewcommand{\arraystretch}{1.15}
\begin{tabular}{@{}p{2.05cm}p{1.35cm}p{1.7cm}p{2.1cm}@{}}
\toprule
\textbf{Característica} & \textbf{Se computa en} & \textbf{Valores} & \textbf{Si no hay evidencia} \\
\midrule
Protección de pago declarada & Cliente & declarada / no declarada & \textit{no verificable} \\
Veracidad de dominio & Backend & escala ordinal & \textit{no verificable} \\
Reputación del vendedor & Backend & escala ordinal & \textit{no verificable} \\
Autenticidad de reseñas & Backend & proporción estimada & \textit{no verificable} \\
Puntuación del producto & Cliente & continua & ausente \\
Presencia de reseñas & Cliente & binaria & ausente \\
Políticas de garantía & Cliente & ordinal & \textit{no verificable} \\
Volumen de ventas & Cliente & entero & ausente \\
\bottomrule
\end{tabular}
\end{table}

La \emph{protección de pago declarada} fija el criterio que rige al resto. Plataformas como
MercadoLibre, Amazon o AliExpress declaran en la ficha que retienen el pago o que garantizan
la devolución ante incumplimiento, y la extensión detecta esa declaración; cuando no la
encuentra, el sistema no infiere que el pago sea inseguro sino que registra la característica
como \emph{no verificable}, según la distinción establecida en la Sección~III-D.

Esta regla se aplica de manera uniforme: para toda característica se distingue entre valor
conocido, valor ausente por no aplicar y valor indeterminado por falta de evidencia. Una
característica se considera \emph{resuelta} cuando toma un valor conocido; los otros dos
estados la dejan sin resolver, y la Tabla~\ref{tab:senales} indica cuál de ellos adopta cada
característica cuando la ficha no aporta evidencia. El agregador propaga esa distinción hasta
el veredicto mediante dos condiciones que deben cumplirse conjuntamente.

La primera es de cobertura. Sea $r$ el número de características resueltas de las ocho; el
sistema exige $r \geq 4$. La segunda es de suficiencia por dimensión. Las ocho características
cubren tres dimensiones de riesgo distintas: la identidad del sitio (veracidad de dominio), la
fiabilidad del vendedor (reputación, volumen de ventas, políticas de garantía y protección de
pago declarada) y la evidencia de la comunidad (puntuación, presencia y autenticidad de
reseñas). El sistema exige que al menos una característica de cada dimensión se haya resuelto.

Cuando ambas condiciones se satisfacen, el veredicto se emite con un nivel de confianza
derivado de la cobertura: alta si $r \geq 6$, media en caso contrario. Cuando alguna falla, el
sistema no emite riesgo sino \emph{evidencia insuficiente}, junto con la lista de
características que no pudo resolver.

La segunda condición distingue esta regla de un simple conteo: un umbral global trataría por
igual la pérdida de tres características redundantes y la de la única que informa sobre la
identidad del sitio, que deja al comprador sin protección frente a la suplantación. La
justificación de fondo es la asimetría de costos, pues presentar como seguro aquello que solo
es desconocido traslada al comprador un riesgo que el sistema no ha evaluado.

\subsection{Descomposición en servicios}

El motor de riesgo de la Fig.~\ref{fig:arquitectura} no es un bloque único: se despliega en
servicios desacoplados. Un servicio central expone la interfaz que consumen los clientes y
delega la estimación de riesgo del vendedor y el procesamiento de reseñas en servicios
independientes, cada uno con su propio ciclo de despliegue; la persistencia se resuelve sobre
una base relacional y la identidad mediante un proveedor externo, que permite asociar cada
consulta a su origen sin custodiar credenciales propias, si bien el uso de la extensión no
exige que el comprador disponga de cuenta. La separación responde a una propiedad que el
sistema necesita: el clasificador de riesgo se reentrena cuando cambia el conjunto de
características disponibles y el detector de reseñas cuando cambian los patrones lingüísticos
del fraude, de modo que aislarlos permite sustituir cualquiera de los dos sin interrumpir el
servicio ni versionar la interfaz que consume la extensión.

\subsection{Analizador de veracidad de dominio}

El analizador de dominio recibe el nombre de host de la ficha y devuelve una estimación
ordinal de veracidad. Su entrada se compone de atributos del registro del dominio, de la
cadena de resolución (incluidos los registros de nombre canónico (CNAME) y el espacio de nombres al que delega) y
de la coincidencia con patrones morfológicos asociados a campañas de suplantación.

Este módulo es necesario porque, como se expuso en la Sección~II-D, ni el certificado SSL ni
las listas de bloqueo del navegador discriminan ya la suplantación \cite{ref25}. La ventaja de
analizar registro y resolución en lugar de contenido es que estos atributos son costosos de
falsificar y están disponibles antes de que el usuario interactúe con la página.
Se ha mostrado que características derivadas del registro permiten detectar URL de
\textit{phishing} en dominios aparcados en una institución financiera \cite{ref7}, lo que
sostiene la viabilidad del enfoque sobre esta clase de atributos. El
módulo no reemplaza la comprobación de certificado, que la literatura de seguridad en
comercio electrónico caracteriza como el control fundacional y de menor costo
\cite{ref30}, sino que la complementa en la dimensión donde aquella dejó de ser informativa.

\subsection{Modelos de inferencia sobre vendedor y reseñas}

\begin{table}[t]
\caption{Configuración del clasificador de riesgo del vendedor}
\label{tab:rf}
\centering
\footnotesize
\renewcommand{\arraystretch}{1.12}
\begin{tabular}{@{}p{4.3cm}p{3.4cm}@{}}
\toprule
\textbf{Parámetro} & \textbf{Valor} \\
\midrule
Algoritmo                     & Bosque aleatorio \\
Número de árboles             & 100 \\
Variables de entrada          & 7 \\
Clases                        & Riesgo alto / moderado / bajo \\
Registros                     & 10\,000 \\
Distribución de clases        & 3\,333 / 3\,334 / 3\,333 \\
Partición entrenamiento--prueba & 80 / 20 \\
Semilla                       & 42 \\
\bottomrule
\end{tabular}
\end{table}

\begin{table*}[t]
\caption{Configuración del clasificador de reseñas fraudulentas}
\label{tab:modelo}
\centering
\footnotesize
\renewcommand{\arraystretch}{1.12}
\begin{tabular}{@{}p{3.9cm}p{3.6cm}p{4.1cm}p{4.3cm}@{}}
\toprule
\multicolumn{2}{@{}l}{\textbf{Arquitectura}} & \multicolumn{2}{l@{}}{\textbf{Ajuste fino}} \\
\cmidrule(r){1-2}\cmidrule(l){3-4}
Modelo base                  & DistilBERT              & Optimizador                    & AdamW \textit{fused} \\
Cabezal                      & Clasif. de secuencias   & Tasa de aprendizaje            & $2\times10^{-5}$ \\
Capas del codificador        & 6                       & Planificador                   & Lineal, sin calentamiento \\
Cabezas de atención          & 12                      & Decaimiento de pesos           & 0.01 \\
Dimensión de representación  & 768                     & $\beta_1$, $\beta_2$, $\epsilon$ & 0.9,\ 0.999,\ $10^{-8}$ \\
Dimensión \textit{feed-forward} & 3072                 & Recorte de gradiente           & 1.0 \\
Parámetros                   & 66.96 M (FP32)          & Tamaño de lote (entr./eval.)   & 16 / 16 \\
Activación                   & GELU                    & Acumulación de gradiente       & 1 \\
\textit{Dropout} gen./aten./clasif. & 0.1 / 0.1 / 0.2  & Épocas                         & 1 \\
Longitud máx. de secuencia   & 512 \textit{tokens}     & Precisión numérica             & FP32 \\
Tamaño de vocabulario        & 30\,522                 & Evaluación                     & Por época \\
Clases                       & \textsc{real}, \textsc{fake} & Semilla                   & 42 \\
\bottomrule
\end{tabular}
\end{table*}

Dos componentes de aprendizaje automático completan el motor. El primero estima el riesgo del
vendedor a partir del vector de características extraídas y devuelve un nivel ordinal; el
segundo estima la proporción de reseñas presumiblemente fraudulentas en la ficha. Sus
configuraciones se resumen en las Tablas~\ref{tab:rf} y~\ref{tab:modelo}.

Para la estimación de riesgo del vendedor se emplea un bosque aleatorio y no una arquitectura
de mayor capacidad, por tres propiedades que el problema exige. La dimensionalidad: el vector
de entrada contiene ocho características heterogéneas (binarias, categóricas, discretas y
continuas), un régimen en el que los métodos de ensamble por árboles operan sin normalización
ni codificaciones adicionales. La interpretabilidad: el modelo expone la contribución relativa
de cada característica, con la que se construye la explicación que acompaña al veredicto. Y el
costo: la inferencia sobre cien árboles es despreciable frente a la latencia de red de la
consulta.

El clasificador del que se dispone actualmente difiere de esta especificación: fue entrenado
sobre siete entradas, sin la veracidad de dominio que este trabajo incorpora. Añadirla exige
reentrenar el modelo sobre un conjunto que la contenga, y su ausencia es consecuencia de que
este trabajo describe la arquitectura propuesta y no una implementación ya construida.

El segundo componente estima la autenticidad de las reseñas con DistilBERT, un codificador
transformador destilado cuya configuración de ajuste recoge la Tabla~\ref{tab:modelo}. La
elección de una variante destilada responde a una restricción operativa: la inferencia ocurre
entre la consulta de la extensión y la emisión del veredicto, de modo que su costo incide
sobre el tiempo que el comprador espera. Seis capas y 66.96 millones de parámetros reducen ese
costo frente a un codificador completo, y la ejecución en precisión de 32 bits sobre unidad
central de proceso evita exigir aceleración por hardware. La contrapartida es una capacidad de
representación menor, cuya repercusión sobre la exactitud debe cuantificarse en la validación.

El ajuste fino se realizó durante una época con tasa de aprendizaje de $2\times10^{-5}$ y
planificador lineal, valores habituales para adaptar un codificador preentrenado sin alterar
las representaciones aprendidas. La entrada del clasificador no es únicamente el texto de la
reseña: se serializan conjuntamente la categoría del producto, la calificación asignada y el
texto, de modo que el modelo dispone del contexto en el que la reseña fue emitida. Esta
decisión responde a que una misma formulación puede resultar verosímil o anómala según la
calificación que la acompaña. La brevedad del ajuste, una sola época, constituye una limitación que la Sección~V recoge.
Debe señalarse además que este trabajo no evalúa el clasificador de reseñas: la
Tabla~\ref{tab:modelo} documenta su configuración para hacer reproducible el componente, no
para respaldar afirmación alguna sobre su desempeño. Ni el corpus de ajuste ni sus métricas de
prueba forman parte de lo reportado aquí, y su evaluación queda comprendida en el trabajo
futuro de la Sección~VI.

La inferencia sobre reseñas es necesaria porque, según el resultado de equilibrio citado en
la Sección~II-C, la reputación visible es manipulable por construcción \cite{ref29}: tomar la
puntuación agregada como indicador de confianza coincide con el comportamiento que el vendedor
fraudulento busca inducir. La arquitectura incorpora por ello la estimación de autenticidad de
las reseñas como característica independiente de la puntuación, si bien la forma en que ambas
deben combinarse, y en particular si la segunda debe corregir a la primera antes de la
clasificación, constituye una decisión de diseño abierta que la validación deberá resolver.

\subsection{Agregación, presentación y persistencia}

El agregador combina las características resueltas en un veredicto con cuatro desenlaces
posibles: riesgo alto, moderado o bajo cuando la evidencia disponible permite decidir, y
\emph{evidencia insuficiente} cuando no. Los tres primeros se acompañan del nivel de confianza
y de la lista de características que los sustentan; el cuarto, de las características que no
pudieron resolverse. La presentación es deliberadamente explicativa: junto al veredicto se
muestra qué características se verificaron, cuáles quedaron indeterminadas y qué acredita cada
una, conforme al requisito de trazabilidad que la literatura de detección de fraude ha ido
incorporando frente a los clasificadores opacos \cite{ref28}.

Finalmente, cada consulta se registra en una base de datos junto con las características extraídas y
el veredicto emitido. El registro cumple dos funciones: habilita la trazabilidad de cada
advertencia mostrada al usuario y constituye el conjunto de datos con el que los modelos se
reentrenarán en iteraciones posteriores. Dado que la información asociada a una consulta puede
revelar hábitos de compra, la política de retención y el nivel de anonimización forman parte
del diseño; la Sección~V-B fija los términos en que ambos rigen durante la validación.
